% Part: first-order-logic
% Chapter: natural-deduction
% Section: proving-things-quant

\documentclass[../../../include/open-logic-section]{subfiles}

\begin{document}

\olfileid{fol}{ntd}{prq}

\olsection{\usetoken{P}{derivation} مع المكممات}

\begin{ex}
عند التعامل مع المكممات، يجب أن نتأكد من عدم انتهاك شرط المتغير
المميَّز، ويتطلب ذلك أحيانًا أن نغيّر ترتيب تنفيذ بعض الاستدلالات.
وبوجه عام، من المفيد أن نحاول معالجة القواعد الخاضعة لشرط المتغير
المميَّز أولًا (إذ ستكون في موضع أدنى من البرهان المكتمل).

لنر كيف نعطي !!a{derivation} لل!!{formula}
$\lexists[x][\lnot !A(x)] \lif \lnot \lforall[x][!A(x)]$.
نبدأ كالمعتاد، فنكتب:
\begin{prooftree}
\AxiomC{}
\UnaryInfC{$\lexists[x][\lnot !A(x)]\lif \lnot \lforall[x][!A(x)]$}
\end{prooftree}
نبدأ بكتابة ما يلزم لتسويغ الخطوة الأخيرة باستعمال قاعدة~\Intro{\lif}.
\begin{prooftree}
\AxiomC{$\Discharge{\lexists[x][\lnot !A(x)]}{1}$}
\DeduceC{$\lnot \lforall[x][!A(x)]$}
\DischargeRule{\Intro{\lif}}{1}
\UnaryInfC{$\lexists[x][\lnot !A(x)]\lif \lnot \lforall[x][!A(x)]$}
\end{prooftree}
لما لم تكن هناك قاعدة واضحة نطبقها على~$\lnot \lforall[x][!A(x)]$،
سنرتب ال!!{derivation} بحيث يمكننا استعمال قاعدة~\Elim{\lexists}.
ويجب أن ننتبه هنا إلى شرط المتغير المميَّز، ونختار رمزًا ثابتًا لا يرد
في~$\lexists[x][\lnot !A(x)]$ ولا في أي افتراض يعتمد عليه.
(لكن لما لم ترد أي !!{constant}s، صلح أي اختيار.)
\begin{prooftree}
\AxiomC{$\Discharge{\lexists[x][\lnot !A(x)]}{1}$}
\AxiomC{$\Discharge{\lnot !A(a)}{2}$}
\DeduceC{$\lnot \lforall[x][!A(x)]$}
\DischargeRule{\Elim{\lexists}}{2}
\BinaryInfC{$\lnot \lforall[x][!A(x)]$}
\DischargeRule{\Intro{\lif}}{1}
\UnaryInfC{$\lexists[x][\lnot !A(x)] \lif \lnot \lforall[x][!A(x)]$}
\end{prooftree}
لكي نشتق~$\lnot \lforall[x][!A(x)]$، سنحاول استعمال قاعدة
\Intro{\lnot}: وهذا يقتضي أن نشتق تناقضًا، وربما باستعمال
$\lforall[x][!A(x)]$ افتراضًا إضافيًا. وبالطبع، قد يتضمن هذا التناقض
الافتراض~$\lnot !A(a)$ الذي سيسقطه استدلال~\Elim{\lexists}. ويمكننا
ترتيب ذلك هكذا:
\begin{prooftree}
\AxiomC{$\Discharge{\lexists[x][\lnot !A(x)]}{1}$}
\AxiomC{$\Discharge{\lnot !A(a)}{2}, \Discharge{\lforall[x][!A(x)]}{3}$}
\DeduceC{$\lfalse$}
\DischargeRule{\Intro{\lnot}}{3}
\UnaryInfC{$\lnot \lforall[x][!A(x)]$}
\DischargeRule{\Elim{\lexists}}{2}
\BinaryInfC{$\lnot \lforall[x][!A(x)]$}
\DischargeRule{\Intro{\lif}}{1}
\UnaryInfC{$\lexists[x][\lnot !A(x)]\lif \lnot \lforall[x][!A(x)]$}
\end{prooftree}
يبدو أننا اقتربنا من الحصول على تناقض. وأسهل قاعدة يمكن تطبيقها هي
\Elim{\lforall}، وهي لا تخضع لأي شرط على المتغير المميَّز. ولما كان
يجوز لنا استعمال أي حد نريده محل~$x$ المكمم كليًا، فمن الأنسب أن نواصل
استعمال~$a$ حتى نصل إلى تناقض.
\begin{prooftree}
\AxiomC{$\Discharge{\lexists[x][\lnot !A(x)]}{1}$}
\AxiomC{$\Discharge{\lnot !A(a)}{2}$}
\AxiomC{$\Discharge{\lforall[x][!A(x)]}{3}$}
\RightLabel{\Elim{\lforall}}
\UnaryInfC{$!A(a)$}
\RightLabel{\Elim{\lnot}}
\BinaryInfC{$\lfalse$}
\DischargeRule{\Intro{\lnot}}{3}
\UnaryInfC{$\lnot \lforall[x][!A(x)]$}
\DischargeRule{\Elim{\lexists}}{2}
\BinaryInfC{$\lnot \lforall[x][!A(x)]$}
\DischargeRule{\Intro{\lif}}{1}
\UnaryInfC{$\lexists[x][\lnot !A(x)]\lif \lnot \lforall[x][!A(x)]$}
\end{prooftree}

من المهم، ولا سيما عند التعامل مع المكممات، أن نراجع في هذه المرحلة
أن شرط المتغير المميَّز لم يُنتهك. ولما كانت القاعدة الوحيدة التي
طبقناها والخاضعة لشرط المتغير المميَّز هي~\Elim{\exists}، وكان
المتغير المميَّز~$a$ لا يرد في أي افتراض تعتمد عليه، فهذا
!!{derivation} صحيح.
\end{ex}

\begin{ex}
قد نشتق أحيانًا !!a{formula} من !!{formula}s أخرى. وفي هذه الحالات، قد
تكون لدينا افتراضات غير مسقطة. ومن المهم أن نضبط افتراضاتنا إلى جانب
الغاية النهائية.

لنر كيف نعطي !!a{derivation} لل!!{formula}~$\lexists[x][!C(x,b)]$
من الافتراضين~$\lexists[x][(!A(x) \land !B(x))]$ و
$\lforall[x][(!B(x) \lif !C(x,b))]$. نبدأ كالمعتاد بكتابة النتيجة في
الأسفل.
\begin{prooftree}
\AxiomC{}
\UnaryInfC{$\lexists[x][!C(x,b)]$}
\end{prooftree}

لدينا مقدمتان نعمل بهما. ولاستعمال الأولى، أي لمحاولة إيجاد
!!a{derivation} لـ~$\lexists[x][!C(x, b)]$ من
$\lexists[x][(!A(x) \land !B(x))]$، نستعمل قاعدة~\Elim{\lexists}.
ولما كانت تخضع لشرط المتغير المميَّز، فسنطبق تلك القاعدة أولًا. فنحصل
على الآتي:
\begin{prooftree}
\AxiomC{$\lexists[x][(!A(x) \land !B(x))]$}
\AxiomC{$\Discharge{!A(a) \land !B(a)}{1}$}
\DeduceC{$\lexists[x][!C(x,b)]$}
\DischargeRule{\Elim{\lexists}}{1}
\BinaryInfC{$\lexists[x][!C(x,b)]$}
\end{prooftree}
يشترك الافتراضان اللذان نعمل بهما في~$!B$. وقد يكون من المفيد في هذه
المرحلة أن نطبق~\Elim{\land} لفصل~$!B(a)$.
\begin{prooftree}
\AxiomC{$\lexists[x][(!A(x) \land !B(x)])$}
\AxiomC{$\Discharge{!A(a) 
\land !B(a)}{1}$}
\RightLabel{\Elim{\land}}
\UnaryInfC{$!B(a)$}
\DeduceC{$\lexists[x][!C(x,b)]$}
\DischargeRule{\Elim{\lexists}}{1}
\BinaryInfC{$\lexists[x][!C(x,b)]$}
\end{prooftree}

الافتراض الثاني الذي علينا العمل به هو~$\lforall[x][(!B(x) \lif
!C(x,b))]$. ولما لم يكن هناك شرط على المتغير المميَّز، يمكننا أن
نعوّض عن~$x$ بال!!{constant}~$a$ باستعمال~\Elim{\lforall} لنحصل على
$!B(a) \lif !C(a, b)$. ولدينا الآن كل من~$!B(a) \lif !C(a,b)$
و$!B(a)$. وينبغي أن تكون خطوتنا التالية تطبيقًا مباشرًا لقاعدة
\Elim{\lif}.
\begin{prooftree}
\AxiomC{$\lexists[x][(!A(x) \land !B(x))]$}
\AxiomC{$\lforall[x][(!B(x) \lif !C(x,b))]$}
\RightLabel{\Elim{\lforall}}
\UnaryInfC{$!B(a) \lif !C(a,b)$}
\AxiomC{$\Discharge{!A(a) 
\land !B(a)}{1}$}
\RightLabel{\Elim{\land}}
\UnaryInfC{$!B(a)$}
\RightLabel{\Elim{\lif}}
\BinaryInfC{$!C(a,b)$}
\DeduceC{$\lexists[x][!C(x,b)]$}
\DischargeRule{\Elim{\lexists}}{1}
\insertBetweenHyps{\hspace{-5em}}
\BinaryInfC{$\lexists[x][!C(x,b)]$}
\end{prooftree}
لقد أوشكنا!{} فبتطبيق واحد لـ~\Intro{\lexists} نبلغ غايتنا.
\begin{prooftree}
\AxiomC{$\lexists[x][(!A(x) \land !B(x))]$}
\AxiomC{$\lforall[x][(!B(x) \lif !C(x,b))]$}
\RightLabel{\Elim{\lforall}}
\UnaryInfC{$!B(a) \lif !C(a,b)$}
\AxiomC{$\Discharge{!A(a) 
\land !B(a)}{1}$}
\RightLabel{\Elim{\land}}
\UnaryInfC{$!B(a)$}
\RightLabel{\Elim{\lif}}
\BinaryInfC{$!C(a,b)$}
\RightLabel{\Intro{\lexists}}
\UnaryInfC{$\lexists[x][!C(x,b)]$}
\DischargeRule{\Elim{\lexists}}{1}
\insertBetweenHyps{\hspace{-5em}}
\BinaryInfC{$\lexists[x][!C(x,b)]$}
\end{prooftree}
وبما أننا حرصنا في كل خطوة على ألا تُنتهك شروط المتغير المميَّز،
يمكننا الوثوق بأن هذا !!{derivation} صحيح.
\end{ex}

\begin{ex}
أعط !!a{derivation} لل!!{formula}~$\lnot\lforall[x][!A(x)]$ من
الافتراضين~$\lforall[x][!A(x)] \lif \lexists[y][!B(y)]$ و
$\lnot\lexists[y][!B(y)]$. نبدأ كالمعتاد بكتابة ال!!{formula}
المطلوبة في الأسفل.
\begin{prooftree}
\AxiomC{}
\UnaryInfC{$\lnot\lforall[x][!A(x)]$}
\end{prooftree}
السطر الأخير في ال!!{derivation} نفي، ولذلك فلنجرب استعمال
\Intro{\lnot}. وسيتطلب هذا أن نعرف كيف !!{derive} تناقضًا.
\begin{prooftree}
\AxiomC{$\Discharge{\lforall[x][!A(x)]}{1}$}
\DeduceC{$\lfalse$}
\DischargeRule{\Intro{\lnot}}{1}
\UnaryInfC{$\lnot\lforall[x][!A(x)]$}
\end{prooftree}
كل شيء على ما يرام حتى الآن. يمكننا استعمال~\Elim{\lforall}، لكن
ليس من الواضح هل سيساعدنا ذلك في بلوغ غايتنا. وبدلًا من ذلك، فلنستعمل
أحد افتراضاتنا. إن~$\lforall[x][!A(x)] \lif \lexists[y][!B(y)]$
مع~$\lforall[x][!A(x)]$ يتيحان لنا استعمال قاعدة~\Elim{\lif}.
\begin{prooftree}
\AxiomC{$\lforall[x][!A(x)] \lif \lexists[y][!B(y)]$}
\AxiomC{$\Discharge{\lforall[x][!A(x)]}{1}$}
\RightLabel{\Elim{\lif}}
\BinaryInfC{$\lexists[y][!B(y)]$}
\DeduceC{$\lfalse$}
\DischargeRule{\Intro{\lnot}}{1}
\UnaryInfC{$\lnot\lforall[x][!A(x)]$}
\end{prooftree}
بقي لدينا الآن افتراض أخير نعمل به، ويبدو أن هذا سيساعدنا على بلوغ
تناقض باستعمال~\Elim{\lnot}.
\begin{prooftree}
\AxiomC{$\lnot\lexists[y][!B(y)]$}
\AxiomC{$\lforall[x][!A(x)] \lif \lexists[y][!B(y)]$}
\AxiomC{$\Discharge{\lforall[x][!A(x)]}{1}$}
\RightLabel{\Elim{\lif}}
\BinaryInfC{$\lexists[y][!B(y)]$}
\RightLabel{\Elim{\lnot}}
\BinaryInfC{$\lfalse$}
\DischargeRule{\Intro{\lnot}}{1}
\UnaryInfC{$\lnot\lforall[x][!A(x)]$}
\end{prooftree}
\end{ex}

\begin{prob}
أعط !!{derivation}s تبين ما يأتي:
\begin{enumerate}
\item $\Proves (\lforall[x][!A(x)] \land \lforall[y][!B(y)]) \lif
\lforall[z][(!A(z) \land !B(z))]$.
\item $\Proves (\lexists[x][!A(x)] \lor \lexists[y][!B(y)]) \lif
\lexists[z][(!A(z) \lor !B(z))]$.
\item $\lforall[x][(!A(x) \lif !B)] \Proves \lexists[y][!A(y)] \lif !B$.
\item $\lforall[x][\lnot !A(x)] \Proves \lnot\lexists[x][!A(x)]$.
\item $\Proves \lnot\lexists[x][!A(x)] \lif \lforall[x][\lnot !A(x)]$.
\item $\Proves \lnot\lexists[x][\lforall[y][((!A(x,y) \lif \lnot
!A(y,y)) \land (\lnot !A(y,y) \lif !A(x,y)))]]$.
\end{enumerate}
\end{prob}

\begin{prob}
أعط !!{derivation}s تبين ما يأتي:
\begin{enumerate}
\item $\Proves \lnot\lforall[x][!A(x)] \lif \lexists[x][\lnot!A(x)]$.
\item $(\lforall[x][!A(x)] \lif !B) \Proves \lexists[y][(!A(y) \lif !B)]$.
\item $\Proves \lexists[x][(!A(x) \lif \lforall[y][!A(y)])]$.
\end{enumerate}
(تتطلب هذه كلها قاعدة~$\FalseCl$.)
\end{prob}

\end{document}
